

{% Q535634[24B]

\needspace{21\baselineskip} 
\item \rtask \ponto{\pt} 
Considere o código fonte Python a seguir.

\begin{lstlisting}[style=Python]
def calcular(n):
    resultado = []
    a, b = 0, 1
    while a < n:
        ...I... 
    return resultado
res = calcular(100)
print(res)
\end{lstlisting}

Para que seja exibido [0, 1, 1, 2, 3, 5, 8, 13, 21, 34, 55, 89] a lacuna I precisa ser preenchida corretamente com:

\begin{answerlist}[label={\texttt{\Alph*}.},leftmargin=*]
\ti \begin{lstlisting}[style=Python,numbers=none]
resultado.insert(a)
a, b = b, a+b
\end{lstlisting}

\di \begin{lstlisting}[style=Python,numbers=none]
resultado.append(a)
a, b = b, a+b
\end{lstlisting}

\ti \begin{lstlisting}[style=Python,numbers=none]
resultado.add(a)
a, b = a, a+b
\end{lstlisting}

\ti \begin{lstlisting}[style=Python,numbers=none]
resultado.append(a)
a, b = a+b, b
\end{lstlisting}

\ti \begin{lstlisting}[style=Python,numbers=none]
resultado.add(a)
a, b = b, a+b
\end{lstlisting}

\end{answerlist}


% Resposta Correta: B
% 
% A questão pede para completar o código de uma função em Python que deve gerar uma sequência de Fibonacci até um número n. A sequência de Fibonacci é uma série de números onde cada número é a soma dos dois precedentes, começando tipicamente com 0 e 1.
% 
% O código fornece a estrutura básica com um laço while que continua a executar enquanto a variável a for menor que n. Para que a função retorne a sequência de Fibonacci desejada, precisamos adicionar cada número atual da sequência (armazenado em a) a uma lista (chamada resultado) e depois atualizar as variáveis a e b para os próximos números da sequência.
% 
% A alternativa correta é a B, que utiliza o método append para adicionar o elemento atual a ao final da lista resultado. Em seguida, atualiza os valores de a e b de tal maneira que a recebe o valor de b e b recebe a soma de a e b (ou seja, o próximo número da sequência), o que é feito com a expressão a, b = b, a+b. Este é o processo que gera a sequência de Fibonacci.
% 
% Veja a alternativa correta formatada com as tags HTML:
% resultado.append(a)
% a, b = b, a+b
% 
% As outras alternativas estão incorretas por diversos motivos:
% 
%     A alternativa A usa um método que não existe para listas em Python (insert existe, mas requer dois argumentos).
%     A alternativa C usa add, que é um método para conjuntos (set) e não listas, além de atualizar incorretamente as variáveis a e b.
%     A alternativa D também usa append, mas atualiza as variáveis de forma invertida, o que não produziria a sequência de Fibonacci.
%     A alternativa E repete o erro de usar o método add, que não é aplicável para listas.
% 
% Entender o funcionamento do método append e a lógica de atualização das variáveis na sequência de Fibonacci é essencial para resolver essa questão corretamente.
}


